\section{谁负责free：指针、别名与所有权}

\begin{frame}{free：释放的是分配，不是指针变量}
\begin{columns}[c]
\begin{column}{0.46\textwidth}
\begin{Code}*{c}\lstinputlisting{code/free.inc}\end{Code}
\end{column}
\begin{column}{0.50\textwidth}
\centering
\begin{tikzpicture}[>=Stealth,font=\scriptsize]
    \node[font=\bfseries,text=black!70] at (0,1.15) {1. \texttt{free(p);} 之后};
    \node[draw=VSPPrimary,fill=VSPPrimary!6,minimum width=1.5cm,minimum height=0.65cm] (p1) at (-1.3,0.45) {\texttt{p}（旧地址）};
    \node[draw=black!45,dashed,fill=black!5,minimum width=2.2cm,minimum height=0.65cm,text=black!55] (b1) at (1.5,0.45) {已归还Heap（禁访）};
    \draw[->,dashed,VSPPrimary,thick] (p1.east) -- node[above,font=\tiny,text=VSPPrimary] {悬空!} (b1.west);
    \node[font=\bfseries,text=black!70] at (0,-0.35) {2. \texttt{p = NULL;} 之后};
    \node[draw=VSPPrimary,fill=VSPPrimary!6,minimum width=1.5cm,minimum height=0.65cm] (p2) at (-1.3,-1.05) {\texttt{p = NULL}};
    \node[draw=black!40,minimum width=1.2cm,minimum height=0.55cm,text=black!60] (n2) at (1.2,-1.05) {\texttt{NULL}};
    \draw[->,thick,black!55] (p2.east) -- (n2.west);
\end{tikzpicture}
\end{column}
\end{columns}
\medskip
\begin{itemize}
\item \texttt{free}结束该块的分配，但\textbf{不会自动把\texttt{p}改为\texttt{NULL}}。
\item \texttt{free(NULL)}安全无操作；主动置\texttt{p = NULL}可避免误用悬空指针。
\end{itemize}
\end{frame}

\begin{frame}{哪几种地址能传给free？}
\begin{columns}[c]
\begin{column}{0.48\textwidth}
\begin{Code}*{c}\lstinputlisting{code/invalid-free.inc}\end{Code}
\end{column}
\begin{column}{0.48\textwidth}
\centering
\begin{tikzpicture}[>=Stealth,font=\scriptsize]
    \node[font=\bfseries,text=black!75] at (1.65,1.25) {整租一套房，退房必须交整套房钥匙};
    \filldraw[fill=VSPPrimary!12,draw=VSPPrimary] (0,0) rectangle (1.1,0.65);
    \node at (0.55,0.325) {\texttt{p[0]}};
    \filldraw[fill=VSPPrimary!12,draw=VSPPrimary] (1.1,0) rectangle (2.2,0.65);
    \node at (1.65,0.325) {\texttt{p[1]}};
    \filldraw[fill=VSPPrimary!12,draw=VSPPrimary] (2.2,0) rectangle (3.3,0.65);
    \node at (2.75,0.325) {\texttt{p[2]}};
    \node[draw=VSPPrimary,fill=white,inner sep=2pt] (p0) at (0.55,-0.75) {\checkmark\ \texttt{free(p)}（块首地址）};
    \draw[->,thick,VSPPrimary] (p0.north) -- (0.55,0);
    \node[draw=black!50,dashed,fill=black!4,inner sep=2pt,text=VSPPrimary] (p1) at (2.35,-1.55) {$\times$\ \texttt{free(p + 1)}（内部地址）};
    \draw[->,thick,dashed,VSPPrimary] (p1.north) -- (1.65,0);
\end{tikzpicture}
\end{column}
\end{columns}
\medskip
\alert{只能传NULL，或与仍有效的动态分配返回值完全匹配的块首指针。}\par\smallskip\small 局部变量（Stack）、全局变量（Static）、块内部偏移地址（\texttt{p + 1}）和已释放的地址都不能传给\texttt{free}。
\end{frame}

\begin{frame}{两个指针指向同一块分配：p = NULL后，q安全吗？}
\begin{columns}[c]
\begin{column}{0.46\textwidth}
\begin{Code}*{c}\lstinputlisting{code/alias.inc}\end{Code}
\end{column}
\begin{column}{0.50\textwidth}
\centering
\begin{tikzpicture}[>=Stealth,font=\scriptsize]
    \node[draw=VSPPrimary,fill=VSPPrimary!6,minimum width=1.5cm,minimum height=0.65cm] (p) at (0,0.65) {\texttt{p = NULL}};
    \node[draw=VSPPrimary,fill=VSPPrimary!6,minimum width=1.5cm,minimum height=0.65cm] (q) at (0,-0.65) {\texttt{q}（旧地址）};
    \node[draw=black!45,dashed,fill=black!5,minimum width=2.4cm,minimum height=1.35cm,align=center,text=black!55] (x) at (3.6,0) {同一块动态分配\\（已被\texttt{free(p)}释放）};
    \draw[->,thick,dashed,VSPPrimary] (q.east) -- node[below=2pt,font=\tiny,text=VSPPrimary] {悬空!} (x.west);
\end{tikzpicture}
\medskip
\begin{itemize}\footnotesize
\item \texttt{int *q = p;}只复制地址，没有复制对象。
\item \texttt{p = NULL}不会自动更新别名\texttt{q}。
\item \alert{释放次数看分配次数，不看有几个指针。}
\end{itemize}
\end{column}
\end{columns}
\end{frame}

\begin{frame}{Memory leak：指针没了，堆上的分配还在}
\begin{columns}[c]
\begin{column}{0.48\textwidth}
\begin{Code}*{c}\lstinputlisting{code/leak.inc}\end{Code}
\end{column}
\begin{column}{0.48\textwidth}
\centering
\begin{tikzpicture}[>=Stealth,font=\scriptsize]
    \node[draw=black!45,dashed,fill=black!4,minimum width=2.1cm,minimum height=1.1cm,align=center,text=black!55] (stack) at (0,0) {\texttt{f}的栈帧已销毁\\局部指针\texttt{p}消失};
    \node[draw=VSPPrimary,fill=VSPPrimary!12,minimum width=2.3cm,minimum height=1.1cm,align=center] (heap) at (3.9,0) {Heap中的100个\texttt{int}\\仍然占着内存};
    \draw[->,dashed,black!35,thick] (stack.east) -- node[above,text=VSPPrimary,font=\bfseries] {$\times$ 断联} (heap.west);
\end{tikzpicture}
\end{column}
\end{columns}
\medskip
\begin{itemize}
\item \texttt{p}是automatic局部指针，离开函数\texttt{f}后自动销毁。
\item 但Heap里的对象不会自动释放；丢失唯一地址后，程序再也无法\texttt{free}它。
\end{itemize}
\alert{“指针生命周期结束”不等于“所指对象自动释放”。}
\end{frame}

\begin{frame}{回看create\_point：约定谁负责free}
\begin{Code}*{c}\lstinputlisting[firstline=12,lastline=22,firstnumber=12]{code/create-point.c}\end{Code}
\begin{itemize}
    \item \texttt{create\_point}把动态对象交给调用者（转移\textbf{ownership}）；返回\texttt{NULL}表示失败。
    \item \texttt{main}接管返回的指针，使用结束后负责\texttt{free(p)}。
    \item 局部结构体\texttt{new\_point}随作用域自动结束生命周期，不能\texttt{free(\&new\_point)}。
\end{itemize}
\end{frame}

\begin{frame}{链表：每个节点都是一次独立分配}
\begin{Code}*{c}\lstinputlisting{code/list-node.inc}\end{Code}
\begin{center}
\begin{tikzpicture}[>=Stealth,node distance=0.7cm,
every node/.style={draw=VSPPrimary,minimum height=0.8cm,inner sep=7pt}]
\node (head) {head};
\node[right=of head] (a) {20 | next};
\node[right=of a] (b) {10 | next};
\node[right=of b] (nil) {NULL};
\draw[->,thick] (head)--(a);
\draw[->,thick] (a)--(b);
\draw[->,thick] (b)--(nil);
\end{tikzpicture}
\end{center}
\alert{释放head指向的节点，不会自动释放next指向的节点。}

\end{frame}

\begin{frame}{释放链表：为什么必须先保存next，再free？}
\begin{columns}[c]
\begin{column}{0.47\textwidth}
\begin{Code}*{c}\lstinputlisting{code/list-free.inc}\end{Code}
\end{column}
\begin{column}{0.49\textwidth}
\centering
\begin{tikzpicture}[>=Stealth,font=\scriptsize]
    \node[font=\bfseries,text=VSPPrimary] at (1.8,1.35) {$\times$ 错误：先\texttt{free(head)}，桥已经炸断};
    \node[draw=black!45,dashed,fill=black!6,minimum width=1.6cm,minimum height=0.6cm,text=black!50] (wa) at (0.1,0.6) {已\texttt{free}节点};
    \node[draw=VSPPrimary,fill=VSPPrimary!8,minimum width=1.5cm,minimum height=0.6cm] (wb) at (3.5,0.6) {下一节点};
    \draw[->,dashed,VSPPrimary,thick] (wa.east) -- node[above,font=\tiny,text=VSPPrimary] {读\texttt{next} = UAF!} (wb.west);

    \node[font=\bfseries,text=black!75] at (1.8,-0.15) {$\checkmark$ 正确：先用\texttt{next}探路，再拆当前桥};
    \node[draw=VSPPrimary,fill=VSPPrimary!8,minimum width=1.6cm,minimum height=0.6cm] (ca) at (0.1,-0.85) {\texttt{head}当前节点};
    \node[draw=VSPPrimary,fill=VSPPrimary!8,minimum width=1.5cm,minimum height=0.6cm] (cb) at (3.5,-0.85) {下一节点};
    \draw[->,thick,VSPPrimary] (ca.east) -- (cb.west);
    \node[draw=VSPPrimary,fill=white,inner sep=2pt] (nx) at (3.5,-1.7) {1. \texttt{next}先指向它};
    \draw[->,thick,VSPPrimary] (nx.north) -- (cb.south);
    \node[font=\tiny,text=VSPPrimary] at (0.1,-1.5) {2. 安全\texttt{free(head)}};
\end{tikzpicture}
\end{column}
\end{columns}
\medskip
\alert{先保存next指针，再free当前节点；否则读取已释放节点的next就是use-after-free。}
\end{frame}

\begin{frame}{内存错误小结：先问哪条规则被破坏了}
\small
\setlength{\tabcolsep}{5pt}
\renewcommand{\arraystretch}{1.28}
\begin{tabular}{lll}
\textbf{错误类型} & \textbf{破坏了哪条规则} & \textbf{本节典型现场}\\[3pt]
Out-of-bounds（越界） & 下标超出有效分配范围 & 申请\texttt{n}个元素却访问\texttt{a[n]}\\
Uninitialized read & 读取前未写入确定值 & 局部\texttt{int local;}或\texttt{malloc}后直接读\\
Dangling / UAF & 所指对象生命周期已结束 & 返回\texttt{\&point}、先\texttt{free(head)}再读\texttt{next}\\
Double free（重复释放） & 同一块分配被\texttt{free}两次 & 别名\texttt{free(p)}后又执行\texttt{free(q)}\\
Memory leak（内存泄漏） & 丢失地址，遗漏释放路径 & \texttt{p = realloc(p, ...)}失败覆盖原指针
\end{tabular}
\medskip\par\alert{这些问题没有按stack、heap、global划分的固定严重程度。}
\par\smallskip 越界、use-after-free、double free等问题，可结合\textbf{上节课的AddressSanitizer}检查。
\end{frame}
